\bookchapter{Capítulo 26: Fronteiras atuais da tecnologia mobile}

\chapterintro{Depois de falarmos sobre o “pós-vida” dos celulares antigos, voltemos nosso olhar para o presente. 2024-2026 é um período de explosão concentrada da tecnologia mobile: a IA saiu da nuvem e entrou no bolso, a comunicação via satélite desceu dos flagships para os intermediários, as telas dobráveis passaram de “experimentação” para “uso cotidiano”, e a tecnologia de baterias teve um avanço em nível de materiais como não se via há anos. Neste capítulo, faremos um panorama dessas fronteiras tecnológicas em andamento, ajudando você a decidir “se deve trocar de celular agora” e “qual comprar”.}

\booksection{I. Celulares com IA: marketing vazio ou verdadeira revolução?}

A partir de 2024, praticamente todos os fabricantes Android passaram a rotular seus flagships como “celulares com IA”, e a Apple também lançou o Apple Intelligence no iOS 18. Mas, afinal, que tipo de celular pode ser considerado um “celular com IA”? O setor estabeleceu aproximadamente três critérios:

\begin{itemize}
  \item \textbf{Modelo de linguagem local (on-device)}: o celular consegue rodar localmente modelos de linguagem com cerca de 7B de parâmetros (como versões leves do Qwen, Llama, DeepSeek), realizando tarefas como sumarização de texto, tradução e geração de conteúdo sem depender da nuvem. Como verificar: se as funções de IA ainda funcionarem no modo avião, significa que o modelo é on-device.
  \item \textbf{Poder computacional da NPU (TOPS)}: a capacidade de computação da NPU (Unidade de Processamento Neural) é medida em TOPS (trilhões de operações por segundo). Em 2024, os chips flagships (como Snapdragon 8 Gen3, Dimensity 9300) tinham NPU entre 40-70 TOPS. A nova geração de chips de 2025-2026 geralmente ultrapassa 100 TOPS, sendo este o requisito mínimo de hardware para rodar modelos de linguagem on-device.
  \item \textbf{Cenários de AI Agent}: não apenas "perguntas e respostas", mas também a capacidade de executar cadeias de tarefas entre aplicativos — por exemplo, "organize as fotos da reunião de ontem em um álbum e compartilhe com Fulano". O sistema deve ser capaz de acionar automaticamente vários apps, como Galeria, Arquivos e WeChat, para concluir a tarefa. Este é o diferencial central de um celular com IA em relação a um "celular com assistente de IA".
\end{itemize}

\booksection{II. Comparação das estratégias de IA dos fabricantes: quem está levando a sério e quem está apenas rotulando?}

As rotas de IA de cada fabricante são bem distintas. Na hora de comprar, preste atenção a:

\begin{itemize}
  \item \textbf{Apple Intelligence (iOS 18+)}: ênfase na integração com aplicativos do sistema — sumarização de e-mails, priorização de notificações, polimento de texto, edição de fotos com IA na galeria. Predominantemente com modelos on-device; tarefas complexas vão para a nuvem via Private Cloud Compute. Vantagens: proteção de privacidade e fluidez do sistema. Desvantagem: a implementação em chinês é mais lenta, algumas funções sendo gradualmente suportadas apenas em 2025.
  \item \textbf{Google Gemini Nano}: integrado nos Pixel 8/9, suporta "Circle to Search", Magic Compose, sumarização de gravações, etc. O Gemini Nano é um verdadeiro modelo pequeno on-device, com ecossistema open source amigável, sendo a implementação de referência para celulares Android com IA.
  \item \textbf{Huawei Pangu Large Model}: integrado ao assistente Xiaoyi do HarmonyOS, suporta geração de texto on-device, edição de fotos com IA, remoção de objetos com IA (remover pessoas indesejadas), sumarização de chamadas. Com o HarmonyOS NEXT 5.0, a integração em nível de sistema é ainda mais profunda, com capacidades de IA profundamente integradas ao SO.
  \item \textbf{Xiaomi MiLM (Large Model)}: integrado ao HyperOS, suporta retratos com IA, sumarização com IA, pintura com IA, super-resolução. A Xiaomi adota a rota de "colaboração dispositivo-nuvem": modelos leves rodam no dispositivo e tarefas pesadas vão para a nuvem.
  \item \textbf{OPPO (Andes Large Model) e vivo (Blue Heart Large Model)}: profundamente integrados ao ColorOS 15 e OriginOS 5, com foco em remoção com IA, sumarização de documentos com IA, reconhecimento de imagem com IA com perguntas e respostas, busca de fotos em linguagem natural ("encontre as fotos tiradas na praia no ano passado").
\end{itemize}

\textbf{Sugestão de compra}: se você frequentemente precisa de IA para processar documentos de trabalho e organizar informações, priorize modelos com forte poder computacional on-device e funções de IA integradas ao sistema. Se for apenas para uso ocasional de edição de fotos com IA, os modelos intermediários atuais já são suficientes.

\booksection{III. Comunicação via satélite: de “gimmick” de flagship a padrão nos intermediários}

A comunicação via satélite é um dos avanços mais concretos dos últimos dois anos. Em 2022, quando estreou no Huawei Mate 60 e no iPhone 14, era exclusividade de flagships. Em 2025, já chegou à faixa de 2000-3000 yuans.

\begin{itemize}
  \item \textbf{Mensagens curtas BeiDou chegam aos intermediários}: a função de mensagens curtas do satélite BeiDou (capaz de enviar SMS de até 1000 caracteres chineses + localização) foi estendida a modelos intermediários como Redmi, realme e iQOO, sem taxa mensal, uso gratuito. Ideal para trilhas ao ar livre e viagens de carro.
  \item \textbf{Chamadas via satélite Tiantong}: as séries Huawei Mate 60/70 e Honor Magic6 suportam chamadas via satélite Tiantong-1, permitindo fazer/receber chamadas via satélite e enviar mensagens via satélite. É necessário contratar um plano de satélite da operadora (aproximadamente 10-200 yuans/mês). Adequado para quem frequenta áreas sem cobertura.
  \item \textbf{Mensagens via satélite do iPhone (SOS)}: iPhone 14 e posteriores suportam SOS via satélite, permitindo enviar mensagens de texto para serviços de emergência e compartilhar localização. Serviço gratuito por 2 anos; renovação paga disponível em algumas regiões.
  \item \textbf{Sugestão prática}: usuários urbanos basicamente não precisam, mas entusiastas de atividades ao ar livre e profissionais que viajam a regiões remotas são fortemente aconselhados a escolher modelos com função de satélite — pode salvar vidas em momentos críticos.
\end{itemize}

\booksection{IV. Maturidade do ecossistema de telas dobráveis: já podem ser o celular principal?}

As telas dobráveis passaram de “experimentação tecnológica” para “usável como dispositivo principal”. Para avaliar a maturidade, observe três pontos: durabilidade do hardware, adaptação de aplicativos e ecossistema de interação.

\begin{itemize}
  \item \textbf{Durabilidade do hardware}: a vida útil das dobradiças geralmente atinge mais de 500 mil ciclos (100 aberturas/fechamentos por dia dariam 13 anos). A dobradiça em gota d'água reduz os vincos da tela. A certificação IPX8 já é padrão. A tela interna com vidro UTG + película protetora torna a tela mais resistente.
  \item \textbf{Adaptação de aplicativos}: os principais apps (WeChat, Taobao, Douyin, Bilibili, Office, WPS) já foram adaptados para telas grandes, mas cerca de 30\% dos apps de cauda longa ainda não estão adaptados, aparecendo esticados ou com barras pretas. Sugere-se testar seus apps mais usados antes de comprar.
  \item \textbf{Interação no modo suspenso (tenda)}: ao dobrar entre 90° e 120°, alguns apps suportam o modo "conteúdo na metade superior, controles na metade inferior", ideal para videochamadas, selfies e assistir vídeos. Selfies no modo suspenso são uma experiência exclusiva das telas dobráveis.
  \item \textbf{Tela dividida e multitarefa}: na tela grande, é possível usar dois apps lado a lado ou até três apps em janelas flutuantes simultaneamente. A lógica de tela dividida do HyperOS, ColorOS e OriginOS já está bem madura, mas o iOS (estilo iPadOS) ainda não liberou essa função no iPhone.
  \item \textbf{Preço}: os modelos de abertura horizontal (tipo livro) caíram da faixa de 10 mil yuans para 5000-7000 yuans. Os de abertura vertical (tipo flip) já chegaram a 4000-5000 yuans. Em 2025, alguns fabricantes lançaram modelos flip na faixa de 3000 yuans, fazendo as telas dobráveis entrarem de fato no mercado de massa.
\end{itemize}

\textbf{Sugestão de compra}: se você é um usuário intensivo de leitura, processamento de documentos e vídeos, os modelos de abertura horizontal valem a pena. Se você busca portabilidade e atributos sociais, os modelos flip são mais adequados. Recomenda-se passar pelo menos 30 minutos testando em uma loja física antes de decidir.

\booksection{V. 6G e Wi-Fi 7: a velocidade e as fronteiras da conectividade de próxima geração}

Embora o 6G ainda não tenha sido realmente implementado, as tecnologias relacionadas já estão a caminho:

\begin{itemize}
  \item \textbf{Progresso do 6G}: o 3GPP planeja concluir a padronização do 6G em 2027, com comercialização por volta de 2030. A taxa de pico teórica do 6G pode chegar a 1 Tbps (50 vezes o 5G), com latência de até 0,1 ms, suportando "Percepção e Comunicação Integradas" (ISAC). O celular poderia até funcionar como um radar de baixa precisão. Atualmente, os fabricantes ainda estão na fase de validação de protótipos, longe de produtos de consumo.
  \item \textbf{Wi-Fi 7 (802.11be)}: a partir de 2024, tornou-se padrão nos flagships, com taxa de pico teórica de 46 Gbps (4 vezes o Wi-Fi 6), suportando largura de banda de 320 MHz, 4K-QAM, MLO (Multi-Link Operation). Em testes reais com roteador Wi-Fi 7, os flagships alcançam taxas de transferência de 2-5 Gbps, 2-3 vezes mais rápido que o Wi-Fi 6E. Se você trocar para um celular com Wi-Fi 7, lembre-se de adquirir também um roteador Wi-Fi 7 para aproveitar todo o desempenho.
  \item \textbf{Bluetooth 5.4 e LE Audio}: os novos celulares de 2024-2026 geralmente suportam Bluetooth 5.4 e LE Audio (Áudio de Baixa Energia), permitindo conectar dois fones de ouvido simultaneamente para ouvir a mesma música. O Auracast broadcast audio permite que várias pessoas ouçam a mesma fonte de áudio sincronizadamente em locais públicos.
\end{itemize}

\booksection{VI. Baterias de ânodo de silício-carbono: finalmente um avanço material na autonomia}

A autonomia dos celulares foi criticada por anos como “uma carga por dia”. A partir de 2024, a popularização da tecnologia de baterias com ânodo de silício-carbono trouxe uma melhoria substancial.

\begin{itemize}
  \item \textbf{Princípio técnico}: a capacidade teórica do ânodo de grafite tradicional é de cerca de 372 mAh/g. A capacidade teórica do silício chega a 4200 mAh/g (cerca de 10 vezes mais), mas o silício se expande até 300\% em volume durante a carga/descarga, tendendo a pulverizar. O ânodo de silício-carbono, através de uma estrutura composta de nano silício-carbono, controla o problema de expansão do silício dentro de limites aceitáveis, aumentando a densidade de energia da bateria em 15\%-30\%.
  \item \textbf{Desempenho real}: os flagships de 2024-2025 com ânodo de silício-carbono (como Xiaomi 14 Ultra, Honor Magic 6 Ultimate, vivo X100 Pro, etc.) conseguem acomodar baterias de 5000-6000 mAh em corpos de 7-8 mm, cerca de 800-1000 mAh a mais que a geração anterior. Uso intensivo pode durar um dia e meio.
  \item \textbf{Compatibilidade com carregamento rápido de 100W+}: o ânodo de silício-carbono tem boa compatibilidade com carregamento rápido. Carregar 100\% em 17 minutos com 100W/120W já não é novidade. Em 2025, alguns modelos suportam 150W-240W, atingindo 80\% em 10 minutos.
  \item \textbf{Sugestão de compra}: se você é um usuário intensivo, priorize modelos pós-2024 com a combinação "ânodo de silício-carbono + carregamento rápido de 100W+". A ansiedade de bateria será drasticamente reduzida.
\end{itemize}

\booksection{VII. Resumo do capítulo: lembre-se destas fronteiras em uma frase}

\begin{itemize}
  \item Três critérios do celular com IA: modelo de linguagem on-device, poder computacional TOPS da NPU, cenários de AI Agent entre apps.
  \item Comunicação via satélite: mensagens curtas BeiDou chegaram aos intermediários; chamadas Tiantong ainda são exclusivas de flagships.
  \item Telas dobráveis: hardware maduro, adaptação de apps em 70\%, modo suspenso/tela dividida são experiências exclusivas.
  \item 6G ainda está a caminho, Wi-Fi 7 já chegou, Bluetooth LE Audio começa a se popularizar.
  \item Ânodo de silício-carbono + carregamento rápido de 100W+: duplo avanço em autonomia e carregamento.
\end{itemize}

\begin{tcbpractice}{Experimente: faça uma "lista de requisitos técnicos" para seu próximo celular}
Pegue uma folha de papel e, nas categorias “Essencial / Diferencial / Desnecessário”, liste suas necessidades quanto a IA, satélite, tela dobrável, Wi-Fi 7, carregamento rápido e capacidade de bateria. Assim, na próxima vez que for a uma loja física, você conseguirá julgar rapidamente se os “diferenciais” recomendados pelo vendedor são necessidades reais ou apenas tentativa de venda. Por exemplo:

\begin{itemize}
  \item Essencial: bateria de 5000mAh ou mais, carregamento rápido de 80W ou mais, NFC, certificação IP68.
  \item Diferencial: sumarização com IA on-device, Wi-Fi 7, mensagens curtas BeiDou.
  \item Desnecessário: tela dobrável (prefiro o formato tradicional), chamadas via satélite (não viajo para áreas remotas), carregamento rápido extremo de 240W.
\end{itemize}

\end{tcbpractice}

\begin{tcbsummary}{Resumo do capítulo}
\begin{itemize}
  \item Três critérios do celular com IA: modelo on-device, TOPS da NPU, cenários de AI Agent
  \item Estratégias de IA dos fabricantes: Apple foca em privacidade e integração; Google em open source on-device; Huawei/Xiaomi/OPPO/vivo cada um com seus pontos fortes
  \item Comunicação via satélite: mensagens curtas BeiDou nos intermediários; chamadas Tiantong ainda são exclusivas de flagships
  \item Telas dobráveis: hardware maduro, adaptação de apps em cerca de 70\%, modo suspenso e tela dividida são experiências exclusivas
  \item Padrão 6G será definido em 2027; Wi-Fi 7 já se popularizou; Bluetooth LE Audio começa a ser implementado
  \item Baterias de ânodo de silício-carbono tornam 5000-6000mAh o padrão nos flagships, reduzindo drasticamente a ansiedade de bateria
\end{itemize}

\end{tcbsummary}
